\chapter*{Úvod}
\addcontentsline{toc}{chapter}{Úvod}
\markboth{Úvod}{}

Od izolování a identifikování atomu vodíku britským chemikem a fyzikem \textit{Henrym Cavendishem} v roce $1766$, tedy za dob rozvíjející se průmyslové revoluce, uplynulo již necelých 260 let. Od té doby lidstvo zpracovává tento nejrozšířenější prvek ve vesmíru v různých variantách.

Vodík se však v posledních letech dostává do popředí zájmu jako klíčová složka pro budoucí energetické systémy. Díky svým jedinečným vlastnostem je považován za potenciální řešení lokálních i globálních výzev spojených s klimatickou neutralitou a energetickou stabilitou. Vodík může být využit v mnoha průmyslových odvětvích, od chemického průmyslu přes metalurgii až po dopravu. Jeho produkce a využití však čelí řadě technických a ekonomických výzev, které musí být překonány, aby mohl plně nahradit fosilní paliva.

V posledních desetiletích se výrazně zvýšila poptávka po vodíku v rafinériích, což souvisí s zintenzivněním rafinérských procesů a rostoucími nároky na hydrogenaci a hydrokrakování. Předpokládá se, že tento trend bude pokračovat, protože celosvětové standardy paliv nadále snižují přijatelné limity obsahu síry.

Cílem této práce je nejprve obecně popsat a znázornit důležitost vodíku jako potencionálního nosiče energie, vyjmenovat a popsat metody výroby a využití vodíku nejen v rafinérském průmyslu, aby se získal základní přehled o této problematice a kontext pro následující detailnější popis aplikací vodíku v reálných projektech po Evropě. V praktické části je navrhnuto, jak by částečné nahrazení enviromentálně škodlivých metod mohlo být prováděno.